\section{Introduction}\label{sec:introduction}
Protecting every stored neural-network parameter can be more expensive than protecting a selected subset. Vulnerability ranking therefore has a practical role in embedded-system design: it identifies locations whose protection might preserve application behavior under a limited budget. Yet an inexpensive ranking typically measures a different object from the protection decision. The numerical jump of an encoded bit, a gradient-derived score, an exact isolated fault, and the value of suppressing that fault in an already corrupted model are not interchangeable.

This paper studies that distinction with frozen quantized executables. Let $R(A)$ denote the task risk after exactly the stored bits in $A$ have been XOR-flipped. The isolated effect is $R(\{b\})-R(\varnothing)$. Under simultaneous faults, the relevant difference instead involves $R(G\cup\{b\})-R(G)$ for a random background $G$. Both the fault process and the currently protected coordinates determine the distribution of $G$. Even an exact isolated-bit ranking can therefore choose a different protection mask from a context-dependent analysis. This statement follows from elementary counterexamples; whether the discrepancy matters in a real network is an empirical question.

Prior sensitivity-guided fault injection, selective protection, and learning-based allocation already recognize that bits differ in importance and that simple significance rules have exceptions~\cite{weng2024fkeras,weng2026priorifi,catalan2025redundancy,alobaidi2026selectiveeccdataset}. We do not claim to originate that observation, Boolean influence, or adaptive protection. Our focus is narrower: make the exact encoded-bit decision object explicit; separate different sources of interaction-driven prediction error; and measure whether improved conditional description translates into a statistically resolved protection benefit against a strong exact-isolated comparator.

\paragraph{Research questions.}
\textbf{RQ1:} How do structural, gradient, and Hessian-based static rankings relate to exact finite-bit effects and actual protected masks? \textbf{RQ2:} How much does an outside fault background change isolated effects and equal-budget mask orderings? \textbf{RQ3:} Does a restricted empirical selector improve independently evaluated risk over a frozen static allocation and an exact-isolated allocation? \textbf{RQ4:} When does a cheaper outside-conditioned singleton approximation reproduce the exact bounded-candidate decision, and what can be certified about its error?

\paragraph{Contributions.}
First, we specialize established multilinear-risk identities to exact quantized execution, including protection and equal-budget swap identities, nonmonotone counterexamples, and the distinction between weight derivatives and fault-probability derivatives. Second, we provide an exact decomposition of equal-budget contrasts into clean isolated, outside-background, and within-candidate components, with no truncation of candidate interactions in the reference calculation. Third, we report all recorded pools from six campaign/BER arms on three frozen model families, including complete-file anomaly-detection risk. Fourth, we give a bounded-risk remainder and finite-field regret analysis, tested against saved exact tables and separately labeled engineering checks.

The results constrain the narrative as much as they support it. A replication of one historically chosen eight-bit pool finds a small favorable held-out contrast. Subsequent probability-sampled six-bit pools do not resolve selected-versus-exact-isolated superiority in any campaign arm. Outside conditioning explains most observed selection-order changes, but those changes are not themselves a demonstration of better final decisions. This distinction is central to the contribution, not a caveat deferred to future work.

All measured protection is \emph{ideal logical event suppression} on stored INT8 kernels. We neither rank physical storage exposure nor claim measured ECC/TMR, ASIC area, power, timing, radiation reliability, or an optimal whole-network allocation. Sections~\ref{sec:related}--\ref{sec:allocation} establish the objects and methods; Sections~\ref{sec:experiments}--\ref{sec:results} describe the evidence and results; Section~\ref{sec:discussion} assesses their limits. Appendix~\ref{app:math} supplies proofs and the explicitly nonempirical spectral interpretation.
